% =============================================================================
% Capítulo 6 — Resultados
% =============================================================================
\chapter{Resultados}
\label{cap:resultados}

Este capítulo apresenta os resultados quantitativos obtidos na avaliação
comparativa entre os chefes FSM e PPO. Os dados são expostos de forma
descritiva, sem interpretação ou análise de causalidade --- a discussão dos
resultados é realizada no \autoref{cap:discussao}.

% -----------------------------------------------------------------------------
\section{Convergência do Treinamento}
\label{sec:convergencia}

O treinamento do agente PPO foi executado no \textit{run}
\texttt{boss\_v1\_500k}, que atingiu o limite configurado de 500.000 passos de
treinamento em aproximadamente 2.008 segundos (cerca de 33 minutos). A
\autoref{tab:metricas-treino} apresenta as métricas finais registradas pelo
TensorBoard.

\begin{table}[htb]
    \centering
    \caption{Métricas finais do treinamento PPO (\textit{run}
             \texttt{boss\_v1\_500k}).}
    \label{tab:metricas-treino}
    \begin{tabular}{l r}
        \toprule
        \textbf{Métrica}                 & \textbf{Valor} \\
        \midrule
        Recompensa média acumulada       & $1\,766{,}4$ \\
        Desvio-padrão da recompensa      & $12{,}8$ \\
        Tempo total de treinamento       & $2\,007{,}5$ s \\
        Passos totais                    & $500\,000$ \\
        \bottomrule
    \end{tabular}
    \fonte{Dados coletados via TensorBoard.}
\end{table}

O baixo desvio-padrão da recompensa ($12{,}8$) em relação à média
($1\,766{,}4$) indica que, nos episódios finais do treinamento, o agente
apresentava desempenho consistente, conforme ilustrado na evolução da recompensa acumulada (\autoref{fig:cumulative-reward}). Além disso, a convergência do modelo também é evidenciada pela estabilização do erro de valor (\textit{Value Loss}), demonstrando que a rede neural aprendeu a prever as recompensas do ambiente com precisão ao longo dos 500 mil passos (\autoref{fig:tb-loss}).

 \begin{figure}[htb]
     \centering
     \includegraphics[width=0.85\textwidth]{imagens/Environment-Cumulative Reward}
     \caption{Evolução da recompensa acumulada ao longo do treinamento
              (\textit{Environment/Cumulative Reward}).}
     \label{fig:cumulative-reward}
     \fonte{Captura do TensorBoard.}
 \end{figure}

 \begin{figure}[htb]
    \centering
    \includegraphics[width=0.85\textwidth]{imagens/Losses-Value Loss.png}
    \caption{Evolução do erro de valor (\textit{Value Loss}) indicando a estabilização do aprendizado da rede neural.}
    \label{fig:tb-loss}
    \fonte{Captura do TensorBoard.}
\end{figure}

% TODO: Inserir gráfico de Episode Length
% \begin{figure}[htb]
%     \centering
%     \includegraphics[width=0.85\textwidth]{Environment-Episode Length}
%     \caption{Evolução do comprimento dos episódios ao longo do treinamento
%              (\textit{Environment/Episode Length}).}
%     \label{fig:episode-length}
%     \fonte{Captura do TensorBoard.}
% \end{figure}

% -----------------------------------------------------------------------------
\section{Comparativo Quantitativo: FSM vs.\ PPO}
\label{sec:comparativo}

Esta seção apresenta os dados numéricos das 20 sessões oficiais da
avaliação: 10 contra o chefe FSM e 10 contra o chefe PPO.

\subsection{Resultados Gerais}
\label{subsec:resultados-gerais}

A \autoref{tab:resultados-gerais} sintetiza os resultados globais por
condição experimental.

\begin{table}[htb]
    \centering
    \caption{Resultados gerais da avaliação por condição experimental.}
    \label{tab:resultados-gerais}
    \begin{tabular}{l r r r r}
        \toprule
        \textbf{Condição}
            & \textbf{Sessões}
            & \textbf{Vitórias}
            & \textbf{Derrotas}
            & \textbf{Taxa de vitória} \\
        \midrule
        FSM & 10 & 10 & 0 & 100\% \\
        PPO & 10 &  9 & 1 &  90\% \\
        \bottomrule
    \end{tabular}
    \fonte{Dados da coleta oficial.}
\end{table}

Na condição FSM, o jogador venceu todas as 10 sessões. Na condição PPO, o
jogador venceu 9 das 10 sessões, com uma derrota registrada --- a única
ocorrência em toda a amostra de 20 sessões em que o chefe eliminou o jogador.

\subsection{Métricas de Combate}
\label{subsec:metricas-combate}

A \autoref{tab:metricas-comparativas} apresenta as médias e desvios-padrão
das métricas de combate para cada condição.

\begin{table}[htb]
    \centering
    \caption{Médias e desvios-padrão das métricas de combate por condição
             (média $\pm$ desvio-padrão, $n = 10$ por condição).}
    \label{tab:metricas-comparativas}
    \begin{tabular}{l r r}
        \toprule
        \textbf{Métrica}
            & \textbf{FSM}
            & \textbf{PPO} \\
        \midrule
        Duração da luta (s)
            & $5{,}97 \pm 2{,}79$
            & $15{,}18 \pm 7{,}74$ \\
        \addlinespace
        Dano recebido pelo chefe
            & $100{,}00 \pm 0{,}00$
            & $95{,}00 \pm 15{,}81$ \\
        \addlinespace
        Dano recebido pelo jogador
            & $50{,}00 \pm 14{,}14$
            & $56{,}00 \pm 22{,}21$ \\
        \addlinespace
        Pulos do jogador
            & $0{,}30 \pm 0{,}48$
            & $3{,}50 \pm 3{,}44$ \\
        \addlinespace
        Ataques do jogador
            & $10{,}70 \pm 0{,}95$
            & $15{,}80 \pm 2{,}57$ \\
        \addlinespace
        \textit{Dashes} do jogador
            & $0{,}20 \pm 0{,}63$
            & $0{,}50 \pm 1{,}08$ \\
        \addlinespace
        Ataques do chefe
            & $5{,}60 \pm 1{,}58$
            & $296{,}60 \pm 186{,}86$ \\
        \bottomrule
    \end{tabular}
    \fonte{Dados da coleta oficial (20 sessões).}
\end{table}

Observa-se que a duração média das sessões na condição PPO ($15{,}18$ s) é
aproximadamente 2,5 vezes superior à da condição FSM ($5{,}97$ s). O dano
recebido pelo chefe foi de $100{,}00$ em todas as sessões FSM, enquanto na
condição PPO a média foi de $95{,}00$, refletindo a sessão em que o jogador
foi derrotado antes de eliminar o chefe. O dano recebido pelo jogador foi
ligeiramente superior na condição PPO ($56{,}00$) em relação à FSM
($50{,}00$).

A contagem de ataques do jogador foi maior na condição PPO ($15{,}80$) do que
na FSM ($10{,}70$), assim como a contagem de pulos ($3{,}50$ vs.\ $0{,}30$) e
\textit{dashes} ($0{,}50$ vs.\ $0{,}20$).

A diferença mais acentuada entre as duas condições reside na contagem de
ataques do chefe: $5{,}60 \pm 1{,}58$ na condição FSM contra $296{,}60 \pm
186{,}86$ na condição PPO.

\subsection{Distribuição de Ações do Agente PPO}
\label{subsec:distribuicao-acoes}

Para a condição PPO, o sistema de métricas registrou a contagem individual de
cada uma das seis ações discretas do \texttt{BossAgent} ao longo das 10
sessões oficiais. A \autoref{tab:acoes-ppo} apresenta esses dados.

\begin{table}[htb]
    \centering
    \caption{Distribuição das ações discretas do agente PPO nas 10 sessões
             oficiais.}
    \label{tab:acoes-ppo}
    \begin{tabular}{l r r r}
        \toprule
        \textbf{Ação}
            & \textbf{Total}
            & \textbf{Média $\pm$ DP}
            & \textbf{\% aprox.} \\
        \midrule
        \textit{Attack}
            & 2\,966  & $296{,}60 \pm 186{,}86$ & 39,05\% \\
        \textit{Jump}
            & 2\,123  & $212{,}30 \pm 149{,}10$ & 27,95\% \\
        \textit{Dash}
            &    796  & $79{,}60 \pm 40{,}77$   & 10,48\% \\
        \textit{Move Right}
            &    698  & $69{,}80 \pm 31{,}94$   &  9,19\% \\
        \textit{Move Left}
            &    550  & $55{,}00 \pm 31{,}67$   &  7,24\% \\
        \textit{Idle}
            &    463  & $46{,}30 \pm 25{,}38$   &  6,10\% \\
        \midrule
        \textbf{Total}
            & \textbf{7\,596} & --- & 100\% \\
        \bottomrule
    \end{tabular}
    \fonte{Dados da coleta oficial.}
\end{table}

As ações \textit{Attack} e \textit{Jump} correspondem, juntas, a
aproximadamente 67\% de todas as ações executadas pelo agente. As ações de
movimentação (\textit{Move Left} e \textit{Move Right}) somam 16,43\%, e o
\textit{Idle} representa 6,10\% das decisões. As três ações que participam
do mecanismo de \textit{Reward Shaping} --- \textit{Attack}, \textit{Jump} e
\textit{Dash} --- totalizam 77,48\% das ações executadas.
